\section{从通济渠到淮南：运输、税源与军费的改向}
\label{sec:transport-fiscal-transition}

大业元年的洛阳不是一座孤立的新都。东都的官署、仓储和迁置人口，同通济渠的开凿被同时推上日程；一端是关中与关东之间的行政重心，另一端是通向淮河、江都的水路。朝廷想得到的并不只是更快的船，而是一条能让粮、役夫、文书和军队在不同季节汇集的调度线。河道遗存和城址分期能够说明这项工程确有其空间尺度，却不能把正史所载役夫与死者的数字化成一份可核的劳工账。%
\footnote{东都与通济渠见 ST-605-EASTERN-CAPITAL、ST-605-TONGJI-CANAL，依据《隋书·炀帝纪上》《隋书·宇文恺传》《隋书·食货志》及隋唐洛阳、通济渠考古调查。工程起点、修缮与实际运量须分别讨论。}

这条线很快显示出国家能力的两面。它使隋廷能够把东部资源接向中枢，也使征发、粮运与远征叠加到同一些州县和家庭身上。612年辽东战场的补给难题、613年黎阳仓附近的杨玄感起兵，并非由运河直接“造成”；但它们揭出一个更实际的限制：当军队向东北移动，仓储与道路的控制本身就成为政治资源，朝廷不得不在前线、京畿和地方治安之间取舍。隋亡后，后来者继承的不是一套可以自动运转的水网，而是把中央供给系于区域合作的难题。%
\footnote{征辽与杨玄感事见 ST-612-GOGURYEO-INVASION、ST-613-YANG-XUANGAN。《隋书·炀帝纪下》《隋书·东夷传·高丽》《隋书·杨玄感传》与《资治通鉴》可确定大体序列；兵额、粮额和仓储控制的细节不能脱离其战时叙事。}

唐初以长安为中枢，重新面对同一距离。开元九年括户检田，是把逃户和隐漏户重新纳入赋役名册；二十一年裴耀卿整顿江淮漕运，则试图让南方粮食更稳定地抵达关中。名册和粮船在这里互为条件：没有被看见的户口，税源难以预期；没有仓、船、堰闸与沿线官员，已经征得的粮也不能转成京师与边军的供给。制度志喜欢把它们排成改革项目，沿线居民所经历的登记、差役与市价波动却很少留下同样连续的记录。%
\footnote{宇文融括户与裴耀卿漕运见 ST-721-YUWEN-RONG、ST-733-PEI-YAOQING，依据《旧唐书》宇文融、裴耀卿传及食货志、《新唐书·食货志》《通典》，并可与敦煌、吐鲁番户籍文书和运河仓储遗存互相限制。}

安史战争把这种联系从“盛世”的背景推到战场前面。两京失陷后，河北、河南的旧税源和道路不再可靠；758年刘晏整顿盐铁转运，所要修复的是江淮财赋进入军费的通道。780年两税法又放弃以均田、人丁和庸役为理想前提的旧框架，转而在夏秋两季向仍能登记的资产、土地与住户征取。它们不是一前一后的技术升级：盐利、转运和两税使朝廷仍能支付军队，却也让盐场、江淮水路、使职与地方留截成为新的争夺点。%
\footnote{刘晏与两税法见 LT-758-007、LT-780-020。《旧唐书》刘晏、杨炎传及食货志、《新唐书·食货志五》《通典》保存制度线索；刘晏的权限和两税的地区施行并不整齐，不能倒推出统一税率或岁入。}

到885年，河中盐池已足以把田令孜、王重荣、京畿军镇和出奔的皇帝卷进同一场危机。盐在这里不是抽象的“财政收入”，而是养兵、任吏和维持地方军府的可用现金来源。唐亡以后，开封、成都、杭州、金陵和广州各凭河道、港口、腹地或军镇经营自己的兵粮；后周956--958年向淮南推进，也必须先沿淮河作战，再安置降军、修城、重设漕运和税籍。一次受禅或一场胜利改变了谁能发令，不能替代把战场变成可征税、可供给区域的缓慢工作。%
\footnote{盐池危机见 LT-885-067；后周南征和江北接收见 LT-956-113、LT-958-115。分别可据《旧唐书·僖宗本纪》《旧唐书·王重荣传》及《旧五代史·周世宗纪三》《五代会要》卷30〈州郡〉、《资治通鉴》；盐利账目、割让州数和地方接收速度均不能视为已精确复原。}

\paragraph{一条水路带来的得失。}
隋唐五代的运输转型提高了跨区域调度的上限：都城、边军和官署能够从更远的粮区获得支持。它的代价不是一个固定的“民力竭尽”数字，而是把征发、登记、仓储和军费联在同一批地方节点上；节点一旦被军镇、叛军或区域政权控制，中央的命令便失去落点。两税、盐铁与后周的淮南接收保住了国家继续筹饷的能力，也使财政重心更依赖江淮水路和地方中介。
